% Milne, Abelian Varieties, Chapter III sections 1--7 (PDF pp. 85--110).
% Exercises, remarks, asides, examples, and bibliographic prose are omitted.

\section{Overview and definitions}
\begin{theorem}[Jacobian representability when pointed]\label{thm:milne-jacobian-representability-pointed}\dcref{MI:ch3:1.2}
Assume $C(k)\ne\varnothing$. The functor $P_C^0$ of families of invertible sheaves of degree zero on $C$, modulo pullback from the parameter, is represented by an abelian variety $J$, uniquely determined up to unique isomorphism, called the Jacobian variety of $C$.
\end{theorem}
\begin{proof}[Source proof]
For a pointed curve, sufficiently positive twists identify degree-$r$ effective
divisors with an open-and-closed locus in a symmetric power.  The divisor-to-
Picard map has projective-space fibres, and the quotient by this linear
equivalence is represented by the degree-zero component, an abelian variety.
\end{proof}
\begin{proposition}[Universal divisorial correspondence]\label{prop:milne-jacobian-universal-correspondence}\dcref{MI:ch3:1.3}\uses{thm:milne-jacobian-representability-pointed}
For $P\in C(k)$ and $J=\operatorname{Jac}(C)$ there is a divisorial correspondence $M$ on $C\times J$ universal as follows: for every divisorial correspondence $L$ on $C\times T$ from $(C,P)$ to a pointed variety $(T,t)$, whose fibres on $C$ have degree zero, there is a unique pointed regular map $\varphi:T\to J$ with $(1\times\varphi)^*M\simeq L$.
\end{proposition}
\begin{proof}
Apply representability to the family $L$ and use the normalization along
$\{P\}\times T$ to make the representing map pointed.  Uniqueness is the
uniqueness part of the representing property.
\end{proof}
\begin{lemma}[Smoothness of symmetric powers]\label{lem:milne-symmetric-power-smooth}\dcref{MI:ch3:1.5}\uses{prop:milne-symmetric-power}
For a nonsingular curve $C$, the symmetric power $C^{(r)}$ is nonsingular.
\end{lemma}

\begin{proof}

At the image of $(P,\ldots,P)$ the completed local ring is $k[[X_1,\ldots,X_r]]^{S_r}=k[[\sigma_1,\ldots,\sigma_r]]$, which is regular. At a general point the completed local ring is a tensor product of rings of this form, one for each multiplicity in the support, and is again regular.

\end{proof}

\begin{theorem}[Jacobian over an arbitrary field]\label{thm:milne-jacobian-general}\dcref{MI:ch3:1.6}\uses{prop:milne-jacobian-universal-correspondence,prop:milne-descent-varieties}
There is an abelian variety $J/k$ and a morphism of functors $\iota:P_C^0\to J$ such that $\iota:P_C^0(T)\to J(T)$ is an isomorphism whenever $C(T)\ne\varnothing$. The pair $(J,\iota)$ is uniquely determined up to unique isomorphism, and after any extension over which $C$ has a rational point it identifies $P_C^0$ with $J$.
\end{theorem}
\begin{proof}[Source proof]
Choose a finite Galois extension $k'/k$ for which $C(k')$ is nonempty.  The
pointed construction over $k'$ gives an abelian variety $J'$ and identifies
$P^0_{C_{k'}}$ with $J'$.  For each $\sigma\in\operatorname{Gal}(k'/k)$,
uniqueness of the representing object gives an isomorphism
$\sigma J'\simeq J'$, and these isomorphisms satisfy the cocycle condition.
Proposition~\ref{prop:milne-descent-varieties} therefore descends $J'$ to an
abelian variety $J/k$.  If $C(T)$ is nonempty, the pointed representing
isomorphism over $k'$ is compatible with this descent datum, and its invariant
part gives $P^0_C(T)\simeq J(T)$.  The same construction gives the morphism of
functors in general, and uniqueness follows from the representing property
after base change to $k'$.
\end{proof}
\begin{theorem}[Pointed universal property]\label{thm:milne-jacobian-pointed}\dcref{MI:ch3:1.7}\uses{thm:milne-jacobian-general}
If $P\in C(k)$, there is a divisorial correspondence $M_P$ between $(C,P)$ and $J$ such that, for every divisorial correspondence $L$ between $(C,P)$ and a pointed $k$-scheme $(T,t)$, there is a unique pointed morphism $\varphi:T\to J$ satisfying $(1\times\varphi)^*M_P\simeq L$.
\end{theorem}
\begin{proof}[Source proof]
The section $s:T\to C\times T$, $t\mapsto(P,t)$, gives a splitting
$\Pic(C\times T)=q^*\Pic(T)\oplus\ker(s^*)$.  Thus $P_C^0(T)$ is represented
by the subfunctor of degree-zero families whose restriction to
$\{P\}\times T$ is trivial.  Let $M$ be the universal element corresponding to
$1_J$.  Translating $M$ by a unique point of $J$ makes its restriction to
$\{P\}\times J$ trivial; call the resulting family $M_P$.  Its representing
property is precisely the pointed universal property.
\end{proof}
\begin{lemma}[Pointed property from representability]\label{lem:milne-jacobian-pointed-from-representability}\dcref{MI:ch3:1.8}\uses{thm:milne-jacobian-general}
Theorem~\ref{thm:milne-jacobian-general} implies Theorem~\ref{thm:milne-jacobian-pointed}.
\end{lemma}

\begin{proof}

For $P\in C(k)$, the section $s:T\to C\times T$ identifies $P_C^0(T)$ with the subgroup of classes represented by sheaves trivial on $\{P\}\times T$. The universal element $M$ therefore has a unique translate whose restriction to $\{P\}\times J$ is trivial; this translate is $M_P$ and has the stated universal property.

\end{proof}

\begin{proposition}[Descent of varieties and coherent sheaves]\label{prop:milne-descent-varieties}\dcref{MI:ch3:1.13}
For a finite Galois extension $k'/k$ with group $G$, base change gives an equivalence between quasi-projective $k$-varieties and quasi-projective $k'$-varieties with descent datum. It likewise gives an equivalence between coherent sheaves, and a descended sheaf is locally free when its pullback is locally free.
\end{proposition}
\begin{proof}[Source proof]
Descent is effective for affine schemes by taking invariants and then gluing
over a Galois-stable affine cover.  The same construction for modules gives
coherent-sheaf descent; local freeness can be checked after faithfully flat
base change.
\end{proof}
\begin{proposition}[Descent of the Jacobian]\label{prop:milne-jacobian-descent}\dcref{MI:ch3:1.14}\uses{thm:milne-jacobian-general,prop:milne-descent-varieties}
Let $k'/k$ be finite separable. If Theorem~\ref{thm:milne-jacobian-general} holds for $C_{k'}$, then it holds for $C$.
\end{proposition}

\begin{proof}

After enlarging $k'$ assume it is Galois and $C(k')\ne\varnothing$. The universal sheaf on $C_{k'}\times J'$ induces, by uniqueness, isomorphisms $\sigma J'\to J'$ satisfying the cocycle condition. Proposition~\ref{prop:milne-descent-varieties} descends $J'$. The same descent argument for rigid pairs $(L,\alpha)$ shows that the resulting map $P_C^0(T)\to J(T)$ is an isomorphism whenever $C(T)\ne\varnothing$.

\end{proof}


\section{The canonical maps from $C$ to its Jacobian}
\begin{proposition}[Tangent space of the Jacobian]\label{prop:milne-jacobian-tangent}\dcref{MI:ch3:2.1}\uses{thm:milne-jacobian-general}
The tangent space at zero is canonically $T_0J\simeq H^1(C,\mathcal O_C)$; consequently $\dim J$ is the genus of $C$.
\end{proposition}

\begin{proof}

Using $k[\epsilon]/(\epsilon^2)$, the tangent space of $P_C^0$ is the kernel of restriction from $C_\epsilon$. Since $\mathcal O_{C_\epsilon}=\mathcal O_C\oplus\mathcal O_C\epsilon$, the map $a\mapsto1+a\epsilon$ identifies this kernel with $H^1(C,\mathcal O_C)$; representability identifies it with $T_0J$.

\end{proof}

\begin{proposition}[Differentials of the Abel map]\label{prop:milne-abel-differentials}\dcref{MI:ch3:2.2}\uses{prop:milne-jacobian-tangent}
For $P\in C(k)$, $f^P:C\to J$ induces an isomorphism $f^{P*}:H^0(J,\Omega_J^1)\to H^0(C,\Omega_C^1)$.
\end{proposition}
\begin{proof}
The tangent-space identification gives $T_0J\simeq H^1(C,\mathcal O_C)$.
Dualizing and using Serre duality identifies invariant differentials on $J$
with $H^0(C,\Omega_C^1)$, and the Abel map induces this identification.
\end{proof}
\begin{proposition}[The Abel map is a closed immersion]\label{prop:milne-abel-closed-immersion}\dcref{MI:ch3:2.3}\uses{prop:milne-jacobian-tangent,lem:milne-closed-immersion-criterion}
For $g(C)>0$, $f^P:C\to J$ is a closed immersion.
\end{proposition}

\begin{proof}

If $f^P(Q)=f^P(Q')$, then $Q-Q'$ is principal, impossible for distinct points on a positive-genus curve. The tangent map is injective at $P$ because its dual is evaluation $H^0(C,\Omega_C^1)\to\Omega_{C,P}^1$ and the kernel has dimension $g-1$ by Riemann--Roch; translation handles every point. Lemma~\ref{lem:milne-closed-immersion-criterion} now applies.

\end{proof}

\begin{lemma}[Complete tangent-injective map criterion]\label{lem:milne-closed-immersion-criterion}\dcref{MI:ch3:2.4}
If $V$ is complete over an algebraically closed field, $f:V\to W$ is injective on points and injective on every tangent space, then $f$ is a closed immersion.
\end{lemma}

\begin{proof}

The image is closed by completeness, and tangent injectivity together with Nakayama's lemma makes the local maps $\mathcal O_{W,f(Q)}\to\mathcal O_{V,Q}$ surjective.

\end{proof}

\begin{theorem}[Analytic Jacobian]\label{thm:milne-analytic-jacobian}\dcref{MI:ch3:2.5}\uses{thm:milne-jacobi-inversion,thm:milne-abel-theorem}
For $k=\C$, the canonical surjection $e:H^0(C,\Omega_C^1)^\vee\twoheadrightarrow J(\C)$ induces an isomorphism $H^0(C,\Omega_C^1)^\vee/H_1(C,\Z)\simeq J(\C)$ carrying the analytic Abel map to $f^P$.
\end{theorem}

\begin{proof}

Jacobi inversion represents every linear functional by integrals from $P$ to $g$ points. Abel's theorem says precisely that such a functional maps to zero exactly when it is integration over an integral homology cycle. Thus $\ker(e)=H_1(C,\Z)$.

\end{proof}


\section{The symmetric powers of a curve}
\begin{proposition}[Symmetric power of a variety]\label{prop:milne-symmetric-power}\dcref{MI:ch3:3.1}
For every variety $V/k$ there is a variety $V^{(r)}$ and a symmetric map $\pi:V^r\to V^{(r)}$ whose underlying topological space is the quotient by $S_r$, with $U^{(r)}$ affine for every affine open $U\subset V$ and coordinate ring $\Gamma(U^r,\mathcal O)^{S_r}$. Every symmetric morphism factors uniquely through $\pi$; $\pi$ is finite, surjective, and separable.
\end{proposition}

\begin{proof}
In the affine case take $V=\operatorname{Spec}A$ and $V^{(r)}=\operatorname{Spec}(A^{\otimes_k r})^{S_r}$; patch these constructions over affine opens.
\end{proof}

\begin{proposition}[Smoothness of symmetric powers]\label{prop:milne-symmetric-power-smooth-2}\dcref{MI:ch3:3.2}\uses{prop:milne-symmetric-power}
If $C$ is a nonsingular curve, then $C^{(r)}$ is nonsingular.
\end{proposition}

\begin{proof}

At a divisor with multiplicities $r_1,r_2,\ldots$, the completed local ring is $\widehat{\bigotimes}_i k[[X_1,\ldots,X_{r_i}]]^{S_{r_i}}=k[[\sigma_{i,j}]]$, hence regular.

\end{proof}

\begin{definition}[Relative effective Cartier divisor]\label{def:milne-relative-effective-divisor}\dcref{MI:ch3:3.4}
For $\pi:X\to T$, a relative effective Cartier divisor is an effective Cartier divisor whose associated closed subspace is flat over $T$.
\end{definition}
\begin{lemma}[Sum of relative effective divisors]\label{lem:milne-relative-divisor-sum}\dcref{MI:ch3:3.5}\uses{def:milne-relative-effective-divisor}
The sum of two relative effective divisors on $X/T$ is relative effective.
\end{lemma}

\begin{proof}

Affinely, if the local equations are $g,g'$, then $gg'$ is a non-zero-divisor and $0\to R/gR\xrightarrow{g'}R/gg'R\to R/g'R\to0$ proves flatness over the base.

\end{proof}

\begin{proposition}[Base change of relative divisors]\label{prop:milne-relative-divisor-base-change}\dcref{MI:ch3:3.7}\uses{def:milne-relative-effective-divisor}
The pullback of a relative effective divisor under any base change is a relative effective divisor.
\end{proposition}

\begin{proof}
The affine local non-zero-divisor and flatness conditions remain true after tensoring with the new base ring.
\end{proof}

\begin{proposition}[Fibrewise Cartier criterion]\label{prop:milne-relative-divisor-fibrewise}\dcref{MI:ch3:3.8}\uses{def:milne-relative-effective-divisor}
If $D,X$ are flat over $T$ and every fibre $D_t$ is an effective divisor on $X_t$, then $D$ is a relative effective divisor on $X/T$.
\end{proposition}

\begin{proof}

From $0\to I_D\to\mathcal O_X\to\mathcal O_D\to0$, flatness gives fibrewise exactness and $I_D\otimes k(t)\simeq I_{D_t}$. The fibrewise flatness criterion makes $I_D$ flat over $X$; it is coherent of rank one, hence locally principal generated by a non-zero-divisor.

\end{proof}

\begin{corollary}[Finite flat divisors on a smooth curve family]\label{cor:milne-relative-divisor-finite-flat}\dcref{MI:ch3:3.9}\uses{prop:milne-relative-divisor-fibrewise}
For a proper smooth relative curve $C/T$, a closed subspace $D$ is a relative effective divisor iff it is finite and flat over $T$. A section has degree-one image.
\end{corollary}

\begin{proof}
On a curve over a field an effective divisor is finite; apply the fibrewise criterion, properness, and finiteness of a proper quasi-finite morphism.
\end{proof}

\begin{proposition}[Fibrewise containment of divisors]\label{prop:milne-relative-divisor-containment}\dcref{MI:ch3:3.10}\uses{def:milne-relative-effective-divisor}
If relative effective divisors satisfy $D_t\ge D'_t$ for every $t\in T$, then $D\ge D'$.
\end{proposition}

\begin{proof}
Represent $D$ by $(s,L)$. The condition is that $s$ vanish in $L\otimes\mathcal O_{D'}$; its support is closed in $T$ and contains every fibre, hence is all of $T$.
\end{proof}

\begin{proposition}[Split relative divisors]\label{prop:milne-split-relative-divisor}\dcref{MI:ch3:3.11}\uses{prop:milne-relative-divisor-containment}
If the support of a relative effective divisor $D$ on $C/T$ is a union of sections, then $D$ has a unique expression $D=\sum_i n_i s_i(T)$.
\end{proposition}

\begin{proof}
The fibre multiplicities are constant, and Proposition~\ref{prop:milne-relative-divisor-containment} compares the resulting sum with $D$ in both directions.
\end{proof}

\begin{theorem}[Representability of effective divisors]\label{thm:milne-effective-divisor-representability}\dcref{MI:ch3:3.13}\uses{prop:milne-symmetric-power,prop:milne-relative-divisor-base-change}
For every relative effective divisor $D$ on $C\times T/T$ of degree $r$, there is a unique morphism $\varphi:T\to C^{(r)}$ with $D=(1\times\varphi)^*D_{\rm can}$. Thus $C^{(r)}$ represents the functor $\operatorname{Div}_C^r$.
\end{theorem}

\begin{proof}

After a finite flat cover $T'\to T$, the divisor splits and determines the map to $C^{(r)}$ by its sections. The two pullbacks to $T'\times_TT'$ agree because they pull back the same divisor, so descent gives $\varphi$.

\end{proof}


\section{The construction of the Jacobian variety}
\begin{lemma}[A section gives representability]\label{lem:milne-jacobian-section}\dcref{MI:ch3:4.1}\uses{thm:milne-effective-divisor-representability}
For $r>2g$, if the natural map $f:\operatorname{Div}_C^r\to P_C^r$ has a section, then $P_C^r$ is represented by a closed subscheme of $C^{(r)}$.
\end{lemma}

\begin{proof}

The composite of the section with $f$ is represented by a morphism $C^{(r)}\to C^{(r)}$. Its equalizer with the identity is a closed subscheme and represents the locus of divisors selected by the section.

\end{proof}

\begin{proposition}[Local sections of the divisor map]\label{prop:milne-jacobian-local-sections}\dcref{MI:ch3:4.2}\uses{thm:milne-effective-divisor-representability}
For an $(r-g)$-tuple $\mathbf P$ of rational points and $L_{\mathbf P}=L(\sum P_i)$, there is an open $C_{\mathbf P}\subset C^{(r)}$ parametrizing divisors $D$ with $h^0(D-\sum P_i)=1$. The corresponding subfunctor $P_{\mathbf P}\subset P_C^r$ is represented locally, and $\operatorname{Div}_{\mathbf P}^r\to P_{\mathbf P}$ has a section.
\end{proposition}

\begin{proof}[Source proof]
Let $D_{\rm can}$ be the canonical relative effective divisor on
$C\times C^{(r)}$.  Upper semicontinuity applied to
$\mathcal O_C(D_{\rm can}-\sum P_i)$ gives an open subvariety
$C_{\mathbf P}$ on which $h^0(D-\sum P_i)=1$; it is nonempty by choosing the
remaining $g$ points successively so that the resulting divisor has this
property.  If $L$ represents an element of $P_{\mathbf P}(T)$, then
Riemann--Roch gives $h^1(L\otimes L_{\mathbf P}^{-1})=0$ on every fibre.
Consequently $M=q_*(L\otimes L_{\mathbf P}^{-1})$ is invertible and commutes
with base change.  The canonical section of
$L\otimes L_{\mathbf P}^{-1}\otimes q^*M^{-1}$ gives a relative effective
divisor, and hence a section of
$\operatorname{Div}_{\mathbf P}^r\to P_{\mathbf P}$.
\end{proof}

\begin{corollary}[Local Jacobian charts]\label{cor:milne-jacobian-local-charts}\dcref{MI:ch3:4.3}\uses{lem:milne-jacobian-section,prop:milne-jacobian-local-sections}
The functor $P_{\mathbf P}$ is represented by a closed subvariety of $C_{\mathbf P}$.
\end{corollary}

\begin{proof}
Apply Lemma~\ref{lem:milne-jacobian-section} to the section in Proposition~\ref{prop:milne-jacobian-local-sections}.
\end{proof}


\section{The canonical maps from symmetric powers to the Jacobian}
\begin{theorem}[Fibres and tangent map of the Abel map]\label{thm:milne-symmetric-to-jacobian}\dcref{MI:ch3:5.1}\uses{prop:milne-jacobian-tangent,prop:milne-symmetric-power}
For $r\le g$, $f^{(r)}:C^{(r)}\to W^r$ is birational, and $f^{(g)}:C^{(g)}\dashrightarrow J$ is birational. If $D$ is effective and $F$ is its fibre, then $0\to T_DF\to T_D C^{(r)}\to T_{f^{(r)}(D)}J$ is exact.
\end{theorem}

\begin{proof}[Source proof]
By Serre duality, the annihilator of the image of the differential at $D$ is
$H^0(C,\Omega_C^1(-D))$.  Hence the kernel has dimension
$r-g+h^1(D)$, which is $h^0(D)-1$ by Riemann--Roch, the dimension of the
linear system $|D|$ (the fibre through $D$).  Lemma~
\ref{lem:milne-generic-divisor-independence} gives a nonempty open set on
which $h^0(D)=1$.  On this set the fibres are points, so $f^{(r)}$ is
birational for $r\le g$; for $r=g$ its image is $J$.  The dimension computation
also gives the asserted exact tangent sequence at every divisor.
\end{proof}

\begin{lemma}[Generic addition of points]\label{lem:milne-generic-divisor-independence}\dcref{MI:ch3:5.2}
If $h^1(D)>0$, there is a nonempty open subset of $C$ on which adding a point lowers $h^1$ by one and outside which it does not. Consequently, for $r\le g$ there is an open subset of $C^r$ on which $h^0(\sum P_i)=1$.
\end{lemma}

\begin{proof}
Serre duality identifies $H^1(D+Q)^\vee$ with differentials in $H^0(\Omega^1(-D))$ vanishing at $Q$; choose the complement of the common zero locus and iterate.
\end{proof}

\begin{proposition}[Differentials on symmetric powers]\label{prop:milne-symmetric-differentials}\dcref{MI:ch3:5.3}\uses{prop:milne-jacobian-tangent,prop:milne-symmetric-power}
There are canonical isomorphisms $H^0(C,\Omega_C^1)\simeq H^0(C^r,\Omega^1)^{S_r}\simeq H^0(C^{(r)},\Omega^1)$. Under these identifications, $f^{(r)*}:H^0(J,\Omega_J^1)\to H^0(C^{(r)},\Omega^1)$ is an isomorphism, and a differential has a zero at $D$ precisely when its pullback to $C^{(r)}$ vanishes at $D$.
\end{proposition}
\begin{proof}
On $C^r$, invariant differentials are sums of pullbacks from the factors;
descent through the finite symmetric quotient gives the first isomorphisms.
The tangent-space description of the Abel map identifies the pullback of an
invariant differential with this sum, and the local symmetric coordinates show
that vanishing is preserved at the divisor $D$.
\end{proof}
\begin{lemma}[Symmetric differential identity]\label{lem:milne-symmetric-differential-identity}\dcref{MI:ch3:5.4}\uses{prop:milne-symmetric-differentials}
If $\sigma_i$ are the elementary symmetric functions in $X_1,\ldots,X_r$ and $\tau_j=\sum_iX_i^j,dX_i$, then $\sigma_0\tau_0-\sigma_1\tau_1+\cdots+(-1)^m\sigma_m\tau_m=d\sigma_{m+1}$ for $m\le r-1$.
\end{lemma}

\begin{proof}
Write each elementary symmetric function with one variable omitted, use $\sigma_n=\sigma_n^{(i)}+X_i\sigma_{n-1}^{(i)}$, multiply by $(-1)^nX_i^n$, sum over $n$, and then differentiate and sum over $i$.
\end{proof}

\begin{corollary}[Degree of the ordered Abel map]\label{cor:milne-ordered-abel-degree}\dcref{MI:ch3:5.5}\uses{thm:milne-symmetric-to-jacobian,prop:milne-symmetric-power}
For $r\le g$, the map $f^r:C^r\to W^r$ has degree $r!$.
\end{corollary}

\begin{proof}
It is the composite $C^r\to C^{(r)}\xrightarrow{f^{(r)}}W^r$, whose first map has degree $r!$ and second map is birational.
\end{proof}


\section{The Jacobian as Albanese variety; autoduality}
\begin{proposition}[Albanese universal property]\label{prop:milne-jacobian-albanese}\dcref{MI:ch3:6.1}\uses{thm:milne-symmetric-to-jacobian,thm:milne-jacobian-pointed}
For $P\in C(k)$, every map $\varphi:C\to A$ to an abelian variety with $\varphi(P)=0$ factors uniquely as $\varphi=\alpha\circ f^P$ for a homomorphism $\alpha:J\to A$.
\end{proposition}

\begin{proof}
The symmetric map $(P_1,\ldots,P_g)\mapsto\sum_i\varphi(P_i)$ gives a rational map $J\dashrightarrow A$, regular by the extension theorem, and it is a homomorphism since it sends zero to zero. Uniqueness follows because sums of $g$ points of $f^P(C)$ fill $J$.
\end{proof}

\begin{corollary}[Uniqueness of the universal correspondence]\label{cor:milne-jacobian-correspondence-unique}\dcref{MI:ch3:6.2}\uses{prop:milne-jacobian-albanese}
If $N$ is a divisorial correspondence between $(C,P)$ and $J$ with $(1\times f^P)^*N\simeq L_P$, then $N\simeq M_P$.
\end{corollary}

\begin{proof}
The universal property makes $N=(1\times\varphi)^*M_P$ for a map $\varphi:J\to J$. The hypothesis gives $\varphi f^P=f^P$, so Proposition~\ref{prop:milne-jacobian-albanese} gives $\varphi=1$.
\end{proof}

\begin{corollary}[Correspondences and homomorphisms]\label{cor:milne-jacobian-hom-correspondence}\dcref{MI:ch3:6.3}\uses{prop:milne-jacobian-albanese,thm:milne-jacobian-pointed}
For pointed curves $(C_i,P_i)$ with Jacobians $J_i$, there is a bijection between $\operatorname{Hom}_k(J_1,J_2)$ and isomorphism classes of divisorial correspondences between $(C_1,P_1)$ and $(C_2,P_2)$.
\end{corollary}

\begin{proof}
Use Theorem~\ref{thm:milne-jacobian-pointed} to pass from a correspondence to a map $C_1\to J_2$, then Proposition~\ref{prop:milne-jacobian-albanese}; the inverse pulls back $M_{P_2}$ along $1\times(f^{P_1}\circ\alpha)$.
\end{proof}

\begin{proposition}[Albanese factorization without a point]\label{prop:milne-jacobian-albanese-no-point}\dcref{MI:ch3:6.4}\uses{prop:milne-jacobian-albanese}
For every map $\varphi:C\times C\to A$ with $\varphi(\Delta)=0$, there is a unique homomorphism $\alpha:J\to A$ with $\alpha\circ F=\varphi$, where $F(P_1,P_2)=f^P(P_1)-f^P(P_2)$; this $F$ is defined over $k$ even when $C(k)=\varnothing$.
\end{proposition}

\begin{proof}
After a finite Galois extension choose $P$. Rigidity gives $\varphi(a,b)=\varphi_1(a)+\varphi_2(b)$ with both terms zero at $P$; the diagonal condition gives $\varphi_1=-\varphi_2$, and Proposition~\ref{prop:milne-jacobian-albanese} applies. Galois uniqueness descends $\alpha$.

\end{proof}

\begin{theorem}[Canonical principal polarization]\label{thm:milne-jacobian-principal-polarization}\dcref{MI:ch3:6.6}\uses{prop:milne-jacobian-albanese}
If $\Theta=W^{g-1}$ (translated as necessary), then $\varphi_{L(\Theta)}:J\to J^\vee$ is an isomorphism. Hence $J$ has a canonical principal polarization.
\end{theorem}
\begin{proof}
After extending the ground field to an algebraic closure, the standard
identities for the theta divisor give
$\varphi_{L(-\Theta)}=\varphi_{L(\Theta)}$ and
$\varphi_{L(\Theta_a)}=\varphi_{L(\Theta)}$ for every translate $\Theta_a$.
Lemma~\ref{lem:milne-theta-divisor-inverse-image} describes the inverse images
of these translates under the Abel map, and
Lemma~\ref{lem:milne-jacobian-poincare-comparison} identifies the resulting
line bundle with the universal correspondence.  The inverse map constructed
in Lemma~\ref{lem:milne-jacobian-autoduality-inverse} is both a left and a right
inverse, so $\varphi_{L(\Theta)}$ is an isomorphism.  The construction descends
because the theta divisor is canonical up to translation.
\end{proof}
\begin{lemma}[Inverse image of a theta divisor]\label{lem:milne-theta-divisor-inverse-image}\dcref{MI:ch3:6.7}\uses{thm:milne-symmetric-to-jacobian}
On the largest open set where $f^{(g)}$ has zero-dimensional fibres represented by sums of distinct points, $f^{-1}(\Theta_a)$ is the effective divisor representing the unique degree-$g$ divisor $D(a)$ mapping to $a$; its degree is at most $g$.
\end{lemma}

\begin{proof}
Remove from $J$ the image of the locus in $C^{(g)}$ where the fibre of
$f^{(g)}$ has positive dimension, and the images of the loci of divisors with
repeated points.  These images are closed because the relevant maps are
proper, so their complement is the stated dense open set.  Let
$D(a)=P_1+\cdots+P_g$ be the unique divisor over a point of this open set,
with the $P_i$ distinct.  A point $Q$ lies in $f^{-1}(\Theta_a)$ exactly when
there is an effective divisor $E$ of degree $g-1$ with
$f(Q)+f^{(g-1)}(E)=a$.  Equivalently, $Q+E$ is linearly equivalent to
$D(a)$; uniqueness of the degree-$g$ divisor gives
$Q\in\{P_1,\ldots,P_g\}$.  Finally, the map
$C\times\Theta\to J$, $(Q,b)\mapsto f(Q)+b$, is projective and has degree
$g$ on the open set, since it is obtained from the ordered Abel map of degree
$g!$ and the ordered map to $\Theta$ of degree $(g-1)!$.  Its finite fibres
therefore have scheme-theoretic rank at most $g$, proving the assertion about
the Cartier divisor.
\end{proof}

\begin{lemma}[Poincare sheaf comparison]\label{lem:milne-jacobian-poincare-comparison}\dcref{MI:ch3:6.8}\uses{lem:milne-theta-divisor-inverse-image,thm:milne-jacobian-principal-polarization}
For every $a\in J(k)$, if $f^{(g)}(D)=a$, then $f^*L(\Theta_a)\simeq L(D)$. Moreover $(f\times(-1))^*L'(\Theta)\simeq M_P$ on $C\times J$.
\end{lemma}

\begin{proof}
For $a$ in the dense open set of Lemma~\ref{lem:milne-theta-divisor-inverse-image},
the inverse image of $\Theta_a$ is the divisor $D$; hence
$f^*L(\Theta_a)\simeq L(D)$.  Both sides vary in families, so the equality
holds for every $a$ by separatedness.  On $C\times\{a\}$, the pullback of
$L^0(\Theta)$ along $(f,-1)$ is therefore the same as the restriction of
$M_P\otimes p^*L(gP)$; on $\{P\}\times J$ both are normalized.  The
see-saw principle then gives
$(f\times(-1))^*L^0(\Theta)\simeq M_P$ (with the chosen normalization),
which is the Poincare-sheaf identity.
\end{proof}

\begin{lemma}[Autoduality inverse]\label{lem:milne-jacobian-autoduality-inverse}\dcref{MI:ch3:6.9}\uses{lem:milne-jacobian-poincare-comparison,thm:milne-jacobian-principal-polarization}
The homomorphism $f^\vee:J^\vee\to J$ defined by $(f\times1)^*P\simeq(1\times f^\vee)^*M_P$ is inverse to $\varphi_{L(\Theta)}$.
\end{lemma}

\begin{proof}
Let $\psi=\varphi_{L(\Theta)}$.  Using
$(1\times\psi)^*P\simeq L^0(\Theta)$ and the defining identity
$(1\times f^\vee)^*M_P\simeq(f\times1)^*P$, we obtain
$(1\times f^\vee\psi)^*M_P\simeq M_P$.  The pointed universal property of
$M_P$ forces $f^\vee\psi=1_J$.  Pulling back the same identities in the
opposite order gives $\psi f^\vee=1_{J^\vee}$, so the two maps are inverse.
\end{proof}


\section{Weil's construction of the Jacobian variety}
\begin{lemma}[Generic Riemann--Roch divisors]\label{lem:milne-weil-generic-divisors}\dcref{MI:ch3:7.1}\uses{thm:milne-effective-divisor-representability}
For a pointed curve $(C,P)$ of genus $g$, there are open subsets $U,V\subset C^{(g)}\times C^{(g)}$ such that $h^0(D+D'-gP)=1$ on $U$ and $h^0(D'-D+gP)=1$ on $V$.
\end{lemma}

\begin{proof}
Upper semicontinuity makes the loci open; Riemann--Roch and the existence of divisors with the required $h^0$ make them nonempty.

\end{proof}

\begin{proposition}[Birational group law on $C^{(g)}$]\label{prop:milne-weil-birational-group}\dcref{MI:ch3:7.2}\uses{lem:milne-weil-generic-divisors}
There is a unique rational map $m:C^{(g)}\times C^{(g)}\dashrightarrow C^{(g)}$ defined on $U$ with $m(D,D')\sim D+D'-gP$; it makes $C^{(g)}$ a birational group.
\end{proposition}

\begin{proof}[Source proof]
On $U$, the unique section of the line bundle
$\mathcal O_C(D+D'-gP)$ gives an effective divisor linearly equivalent to
$D+D'-gP$.  The relative divisor representability theorem makes this
construction a rational map $m$.  On the common dense open locus where all
expressions are defined, both $m(m(D,D'),D'')$ and $m(D,m(D',D''))$ are the
unique effective divisor in the class $D+D'+D''-2gP$, so associativity holds.
The analogous construction using $V$ gives birational inverses for
$(D,D')\mapsto(D,m(D,D'))$ and $(D,D')\mapsto(D',m(D,D'))$; the divisor $gP$
is the identity.  Thus $C^{(g)}$ is a birational group.
\end{proof}

\begin{theorem}[Birational group completion]\label{thm:milne-weil-birational-group-completion}\dcref{MI:ch3:7.3}\uses{prop:milne-weil-birational-group,prop:milne-descent-varieties}
For every birational group $V/k$ there is a group variety $G/k$ and a birational map $f:V\dashrightarrow G$ with $f(ab)=f(a)f(b)$ wherever defined; $G$ is unique up to unique isomorphism.
\end{theorem}

\begin{proof}[Source proof]
In the case that $V(k)$ is dense in $V$, the construction is obtained by
replacing $V$ by an open subset on which the rational law has the required
regularity and patching copies of that open subset by its translates (Artin,
1986, \S2).  The same construction applies after a finite Galois extension
$k'/k$, so let $G$ and $f:V_{k'}\dashrightarrow G$ be obtained there.  For
$\sigma\in\operatorname{Gal}(k'/k)$, the conjugate $\sigma f$ is a birational
map from $V_{k'}$ to $\sigma G$.  Uniqueness over $k'$ gives a unique
isomorphism $\phi_\sigma:\sigma G\to G$ such that
$\phi_\sigma\circ\sigma f=f$.  For $\sigma,\tau$ the two maps from
$\tau\sigma G$ to $G$ obtained by composing these isomorphisms agree after
composition with the birational map $\tau\sigma f$; uniqueness therefore
gives $\phi_\tau\circ\tau(\phi_\sigma)=\phi_{\sigma\tau}$.  These
isomorphisms are a descent datum, and Proposition~\ref{prop:milne-descent-varieties}
descends $G$ and $f$ to $k$.

\end{proof}

\begin{proposition}[Completeness of Weil's Jacobian]\label{prop:milne-weil-jacobian-complete}\dcref{MI:ch3:7.4}\uses{thm:milne-weil-birational-group-completion}
The group variety $J$ obtained from the birational group $C^{(g)}$ is complete.
\end{proposition}
\begin{proof}[Source-omitted proof]
Milne says that this can be proved using the valuative criterion of
properness, referring to Weil (1948b), Theoreme~16 and following results, but
does not include that argument here.
\end{proof}
\begin{corollary}[Extension of the rational map]\label{cor:milne-weil-map-morphism}\dcref{MI:ch3:7.5}\uses{prop:milne-weil-jacobian-complete}
The rational map $f:C^{(g)}\dashrightarrow J$ is a morphism. Linearly equivalent divisors over any extension field have the same image in $J$.
\end{corollary}

\begin{proof}
Completeness and the extension theorem extend the rational map; a complete linear system is a projective space, and the map is constant on it by the extension property.
\end{proof}

\begin{theorem}[Weil construction represents the Picard functor]\label{thm:milne-weil-jacobian-representability}\dcref{MI:ch3:7.6}\uses{cor:milne-weil-map-morphism,thm:milne-effective-divisor-representability}
There is a canonical isomorphism of functors $P_C^0\simeq J$.
\end{theorem}

\begin{proof}
It suffices to represent $P_C^g$. Locally where $h^0(L)=1$, the canonical section of $L\otimes(q_*L)^{-1}$ gives a divisor and hence a map to $C^{(g)}$ and then $J$. Where $h^0(L)>1$, tensor by a suitable degree-zero line bundle to reach the first case. These local maps agree on overlaps and define the required functorial map.

\end{proof}

% Source-faithful fragments for Milne, Abelian Varieties, PDF pages 117--172.
% Sections III.1--III.7 are already represented in the adjacent chapter source;
% this file records the source statements on the requested pages (III.8--III.13)
% and the complete Chapter IV sections 1--8.
\section{Relative and generalized Jacobians}

\begin{theorem}[Relative Jacobian]
\label{thm:milne-relative-jacobian-source}

\dcref{MI:ch3:8.1}
\uses{thm:milne-weil-jacobian-representability}
Let $C\to S$ be a proper flat curve with integral fibres.  There
is a group scheme $\mathcal J$ with connected fibres and a morphism of functors
$P^0_{C/S}\to\mathcal J$ that is injective and is an isomorphism whenever
$C\times_S T\to T$ has a section.  The construction commutes with base change.
If $C/S$ is smooth, $\mathcal J$ is an abelian scheme, and its fibre at $t$ is
the Jacobian of $C_t$.
\end{theorem}

\begin{proof}[Source-omitted proof]

Relative effective divisors are represented by a Hilbert scheme.  After a
section is chosen, sufficiently positive twists embed the degree-$r$ divisor
functor in a Grassmann scheme; the fibres of the divisor-to-Picard map are
projective-space bundles.  The resulting quotient gives the required group
scheme.  The same construction gives the injectivity and the assertion when a
section exists, and commutes with base change.

\end{proof}


\begin{proposition}[Grassmann representability]
\label{prop:milne-grassmann-representability-source}

\dcref{MI:ch3:8.3}
If $\mathcal E$ is locally free and $\operatorname{Grass}_{\mathcal E}^n(T)$
classifies rank-$n$ locally free quotients of $\mathcal O_T\otimes\mathcal E$,
then this functor is represented by a projective variety $G_{\mathcal E}^n/S$.
\end{proposition}

\begin{proof}[Source-omitted proof]

This is the usual Pluecker construction of the Grassmann scheme, applied
locally to a free presentation and descended under change of frame.

\end{proof}


\begin{proposition}[Divisors in a Grassmannian]
\label{prop:milne-relative-divisors-grassmann-source}

\dcref{MI:ch3:8.4}
\uses{prop:milne-grassmann-representability-source}
For $r>2g-2$ and $m>2g-2+r$, the map sending a relative effective divisor
$D$ to the quotient $q_*\mathcal O_D(m)$ identifies $\operatorname{Div}^r_{C/S}$
with a closed subscheme of the appropriate Grassmann scheme.  The fibres of
$\operatorname{Div}^r_{C/S}\to P^r_{C/S}$ are projective-space bundles.
\end{proposition}

\begin{proof}[Source-omitted proof]

The exact sequence $0\to\mathcal L(-D)(m)\to\mathcal O_C(m)\to
\mathcal O_D(m)\to0$ has vanishing $R^1q_*$ by Serre duality and the bounds
on $m,r$.  Thus $q_*\mathcal O_D(m)$ is locally free of rank $r$ and gives the
Grassmann map.  Conversely its kernel defines the ideal of a closed subscheme;
the construction is inverse on relative divisors.  The closedness and the
projective-bundle description are the standard Hilbert/Grassmann argument.

\end{proof}


\section{Coverings from the Jacobian}
\begin{proposition}[Abelian coverings]
\label{prop:milne-abelian-coverings-source}

\dcref{MI:ch3:9.1}
Let $C/k$ be a complete nonsingular pointed curve with Jacobian $J$, over a
separably closed field.  Pullback along the Abel map induces an isomorphism
$\pi_1(C,P)^{\mathrm{ab}}\simeq\pi_1(J,0)$; hence every connected abelian
etale cover of $C$ is the pullback of a connected etale cover of $J$.
\end{proposition}

\begin{proof}

For $n$ prime to $\operatorname{char}(k)$, Kummer theory identifies continuous
$\mathbb Z/n\mathbb Z$-characters of $\pi_1(V)$ with $\Pic(V)[n]$.  Pullback
induces $\Pic^0(J)\simeq\Pic^0(C)$, proving the assertion for such $n$.
For $n=p=\operatorname{char}(k)$, the analogous character group is
$\ker(1-F:H^1(V,\mathcal O_V)\to H^1(V,\mathcal O_V))$, and
$H^1(J,\mathcal O_J)\simeq H^1(C,\mathcal O_C)$ by the Abel map.  Witt vectors
give the $p^m$ case; passing to inverse limits proves the assertion.

\end{proof}


\begin{lemma}[Kummer characters]
\label{lem:milne-kummer-characters-source}

\dcref{MI:ch3:9.2}
For a complete nonsingular $V$ and $n$ prime to the characteristic,
$\operatorname{Hom}_{\mathrm{cont}}(\pi_1(V,P),\mathbb Z/n\mathbb Z)
\simeq\Pic(V)[n]$.  The map sends an $n$-torsion divisor class to the unramified Kummer
cover obtained by adjoining an $n$th root of a representative function.
\end{lemma}

\begin{proof}

Purity says the Kummer extension is etale exactly when every codimension-one
valuation of the defining divisor is divisible by $n$; conversely an etale
cyclic cover has this divisibility by Kummer theory.

\end{proof}


\begin{lemma}[Picard comparison for the Abel map]
\label{lem:milne-picard-comparison-source}

\dcref{MI:ch3:9.3}
\uses{prop:milne-abelian-coverings-source}
The Abel map induces $\Pic^0(J)\simeq\Pic^0(C)$.
\end{lemma}

\begin{proof}
 This is the Picard comparison obtained from the universal
divisorial correspondence (the argument of III.6).
\end{proof}


\begin{lemma}[Characteristic-$p$ characters]
\label{lem:milne-characteristic-p-characters-source}

\dcref{MI:ch3:9.4}
For $p=\operatorname{char}(k)$,
$\operatorname{Hom}(\pi_1(V,P),\mathbb Z/p\mathbb Z)
\simeq\ker(1-F:H^1(V,\mathcal O_V)\to H^1(V,\mathcal O_V))$.
\end{lemma}
\begin{proof}[Source proof]
Artin--Schreier theory identifies cyclic $p$-covers with the kernel of
$1-F$ on $H^1(V,\mathcal O_V)$.
\end{proof}

\begin{lemma}[Structure sheaf comparison]
\label{lem:milne-structure-cohomology-source}

\dcref{MI:ch3:9.5}
The Abel map gives $H^1(J,\mathcal O_J)\simeq H^1(C,\mathcal O_C)$.
\end{lemma}
\begin{proof}
Both spaces are dual to the spaces of invariant differentials, and the
differentials proposition identifies those spaces via the Abel map.
\end{proof}

\begin{corollary}[Etale cohomology comparison]
\label{cor:milne-etale-cohomology-comparison-source}

\dcref{MI:ch3:9.6}
\uses{prop:milne-abelian-coverings-source}
For every prime $\ell$, $H^1_{\et}(J,\mathbb Z_\ell)\to
H^1_{\et}(C,\mathbb Z_\ell)$ induced by the Abel map is an isomorphism.
\end{corollary}
\begin{proof}
For $\ell\ne\operatorname{char}(k)$ use the Kummer-character description and
$\Pic^0(J)\simeq\Pic^0(C)$.  In characteristic $p$, use the
Artin--Schreier/Witt description together with the structure-sheaf comparison.
\end{proof}

\begin{proposition}[Ramified coverings and generalized Jacobians]
\label{prop:milne-generalized-jacobian-coverings-source}

\dcref{MI:ch3:9.7}
If $C'\to C$ is a finite abelian cover unramified outside $\Sigma$, there is
a modulus $\mathfrak m$ supported on $\Sigma$ and an etale isogeny
$J'\to J_{\mathfrak m}$ whose pullback by the canonical map is
$C'\setminus f^{-1}(\Sigma)$.
\end{proposition}
\begin{proof}[Source-omitted proof]
The local inertia characters determine a modulus supported on $\Sigma$.
Class-field theory for the generalized Jacobian identifies the resulting
abelian cover with an etale isogeny of $J_{\mathfrak m}$, and pullback along
the Abel map recovers $C'$ away from the modulus.
\end{proof}

\begin{proposition}[Parshin's auxiliary cover]
\label{prop:milne-parshin-auxiliary-cover-source}

\dcref{MI:ch3:9.9}
For a genus-$g$ curve over a number field with good reduction outside $S$,
there is an extension $k'/k$ of degree at most $2^{2g}$, unramified outside
$S$, and a cover $C_P\to C_{k'}$ of degree at most
$2^{2g(g-1)+2g+1}$, ramified exactly over a chosen rational point $P$,
with $C_P$ having good reduction outside $S$.
\end{proposition}

\begin{proof}

Pull back multiplication by $2$ on $J$ to obtain an etale cover of degree
$2^{2g}$, then choose a rational point above $P$ after a bounded unramified
extension.  Use the resulting divisor as modulus and pull back the corresponding
generalized-Jacobian cover.  Hurwitz gives the stated genus and degree bounds.

\end{proof}


\begin{conjecture}[Shafarevich]
\label{conj:milne-shafarevich-source}

\dcref{MI:ch3:9.10}
\notready
For fixed $k,g,S$, only finitely many isomorphism classes of genus-$g$ smooth
complete curves over $k$ have good reduction outside $S$.
\end{conjecture}

\begin{theorem}[Shafarevich implies Mordell]
\label{thm:milne-shafarevich-implies-mordell-source}

\dcref{MI:ch3:9.11}
\uses{conj:milne-shafarevich-source,prop:milne-parshin-auxiliary-cover-source}
If Shafarevich's conjecture holds, then every curve of genus at least two over
a number field has finitely many rational points.
\end{theorem}

\begin{proof}

Put all bounded auxiliary fields in one finite extension $K$.  The covers $C_P$
have bounded genus and good reduction outside a fixed set, so Shafarevich gives
finitely many $C_P$.  De Franchis gives finitely many maps from each $C_P$ to
$C$, while ramification exactly at $P$ makes $P\mapsto(C_P,\phi_P)$ injective.

\end{proof}


\section{Abelian varieties as Jacobian quotients}
\begin{theorem}[Jacobian quotient]
\label{thm:milne-abelian-is-jacobian-quotient-source}

\dcref{MI:ch3:10.1}
Every abelian variety over an infinite field is a quotient of a Jacobian.
\end{theorem}

\begin{proof}[Source proof]
For $\dim A>1$, fix a projective embedding.  Bertini gives a hyperplane
section which is smooth and connected, and repeating this produces a smooth
geometrically connected complete-intersection curve $C\subset A$.  Let
$A_1$ be the image of the homomorphism $J_C\to A$ induced by the inclusion.
It is an abelian subvariety.  If $A_1\ne A$, choose an abelian subvariety
$A_2$ such that $A_1\times A_2\to A$ is an isogeny; then $A_1\cap A_2$ is
finite, so $C\cap A_2$ is finite.  Compose the isogeny with multiplication by
an integer $n>1$ on $A_2$.  The inverse image of $C$ is then disconnected,
because its image over $A_2$ is supported on the finite set $C\cap A_2$,
contradicting the geometric connectedness of inverse images of general
hyperplane sections.  Thus $A_1=A$, and $J_C\to A$ is surjective.  For
$\dim A=1$, use that an elliptic curve is its own Jacobian.
\end{proof}


\begin{lemma}[Finite pullback preserves ampleness]
\label{lem:milne-finite-pullback-ample-source}

\dcref{MI:ch3:10.2}
If $\pi:W\to V$ is finite and $\mathcal L$ is ample on complete $V$, then
$\pi^*\mathcal L$ is ample on $W$.
\end{lemma}

\begin{proof}

For coherent $\mathcal F$, affineness and the projection formula give
$H^i(W,\mathcal F\otimes\pi^*\mathcal L^n)=H^i(V,\pi_*\mathcal F\otimes\mathcal L^n)$,
which vanishes for $i>0$ and $n\gg0$.

\end{proof}


\begin{lemma}[Finite fibres of a general projection]
\label{lem:milne-general-projection-fibres-source}

\dcref{MI:ch3:10.3}
If $V$ is nonsingular projective of dimension at least two and $Z$ is a
general hyperplane section in a projective embedding, then inverse images of
$Z$ under a finite map from a nonsingular $W$ are geometrically connected.
\end{lemma}

\begin{proof}

After base change to an algebraic closure, Bertini gives smoothness and
connectedness for general hyperplanes; the connectedness is preserved under
the finite map by the ample pullback criterion.

\end{proof}


\section{Zeta functions of curves}
\begin{theorem}[Point count from Frobenius]
\label{thm:milne-curve-point-count-source}

\dcref{MI:ch3:11.1}
If $C/\mathbb F_q$ is complete nonsingular and the Frobenius polynomial of
$J_C$ is $P(X)=\prod_i(X-a_i)$, then
$\#C(\mathbb F_q)=1+q-\sum_i a_i$ and
$|\#C(\mathbb F_q)-q-1|\le 2gq^{1/2}$; the roots $a_i$ have absolute value
$q^{1/2}$.
\end{theorem}

\begin{proof}

Let $\pi_C$ be Frobenius on $C$ and $\pi_J$ the induced endomorphism of its
Jacobian.  The graph of $\pi_C$ meets the diagonal in exactly
$\#C(\mathbb F_q)$ points, while $\deg(\pi_C)=q$.  Applying the correspondence
trace formula gives
$\#C(\mathbb F_q)=1-\operatorname{Tr}(\pi_J)+q$.  The trace of $\pi_J$ is
the sum of the roots $a_i$ of its characteristic polynomial, and the Weil
bound for those roots gives
$|\#C(\mathbb F_q)-q-1|\le2gq^{1/2}$.

\end{proof}


\begin{proposition}[Correspondence trace formula]
\label{prop:milne-correspondence-trace-source}

\dcref{MI:ch3:11.2}
For an endomorphism $\alpha$ of $C$ and its induced endomorphism of $J$,
the intersection number satisfies $(\Gamma_\alpha\cdot\Delta)=1-\operatorname{Tr}(\alpha)+\deg(\alpha)$.
\end{proposition}

\begin{proof}

Let $L^0(\Theta)$ be the normalized line bundle on $J\times J$ associated
with the theta divisor.  Pulling it back around the diagram formed by
$f\times f$, $1\times\alpha'$, and the diagonal gives a line bundle on $C$
whose degree is
$(\Gamma_\alpha\cdot\Delta)-\deg(\alpha)$.  The same degree computed via
$\alpha+1$ is the intersection of $f(C)$ with
$D_\Theta(\alpha)=(\alpha+1)^*\Theta-\alpha^*\Theta-\Theta$.
The intersection calculation for the theta class and
Lemma~\ref{lem:milne-weil-polynomial-identity-source} identify this with
$-\operatorname{Tr}(\alpha')+1$.  Rearranging gives
$(\Gamma_\alpha\cdot\Delta)=1-\operatorname{Tr}(\alpha')+\deg(\alpha)$.

\end{proof}


\begin{lemma}[Weil polynomial identity]
\label{lem:milne-weil-polynomial-identity-source}

\dcref{MI:ch3:11.3}
For $A/k$ of dimension $g$ and a divisor $H$, with $D_\alpha=(\alpha+1)^*H-\alpha^*H-H$,
\[
\operatorname{Tr}(\alpha)=g\frac{(H^{g-1}\cdot D_\alpha)}{(H^g)}.
\]
\end{lemma}

\begin{proof}

In $\operatorname{NS}(A)$, quadraticity of pullback gives
$(\alpha+n)^*H=n(n-1)H+n(\alpha+1)^*H-(n-1)\alpha^*H$
and hence
$(\alpha+n)^*H=n^2H+nD_\alpha+\alpha^*H$.  Intersecting the $g$th power
with $H^{g-1}$ and using
$\deg(\alpha+n)=((\alpha+n)^*H^g)/(H^g)$, compare the coefficient of
$n^{2g-1}$ with the coefficient of the characteristic polynomial of $\alpha$.
This coefficient is $\operatorname{Tr}(\alpha)$ and yields the displayed
identity.

\end{proof}


\begin{corollary}[Curve zeta function]
\label{cor:milne-curve-zeta-source}

\dcref{MI:ch3:11.4}
\uses{thm:milne-curve-point-count-source}
\[
Z(C,t)=\frac{P(t)}{(1-t)(1-qt)},\qquad
P(t)=\prod_i(1-a_i t).
\]
\end{corollary}
\begin{proof}
For every $m$, the point-count theorem gives
$\#C(\mathbb F_{q^m})=1+q^m-\sum_i a_i^m$.  Insert this in
$\log Z(C,t)=\sum_{m\ge1}\#C(\mathbb F_{q^m})t^m/m$ and use
$-\log(1-x)=\sum_{m\ge1}x^m/m$.
\end{proof}

\section{Torelli's theorem}
\begin{theorem}[Torelli]
\label{thm:milne-torelli-source}

\dcref{MI:ch3:12.1}
For complete nonsingular curves $C,C'$ over an algebraically closed field,
an isomorphism $\beta$ of their canonically principally polarized Jacobians
induces an isomorphism $\alpha:C\simeq C'$ and a sign
$\varepsilon\in\{1,-1\}$ such that
$f'\alpha=\varepsilon\beta f+c$ for some translation $c$.  If $g(C)\ge2$
and $C$ is not hyperelliptic, $\alpha,\varepsilon,c$ are uniquely determined;
if $C$ is hyperelliptic, $\alpha,c$ are uniquely determined for each choice of
$\varepsilon$.
\end{theorem}

\begin{proof}[Source proof]
The principal polarization identifies the theta divisor with a translate of
$W^{g-1}$.  Lemmas~\ref{lem:milne-torelli-theta-symmetry-source}--
\ref{lem:milne-torelli-intersections-source} describe its translates and their
incidence.  Applying these identities successively recovers from the
polarized variety the subvariety $W^1$ (or its inverse $-W^1$).  Since the Abel
map is a closed immersion, $W^1$ recovers $C$ and yields an isomorphism
$\alpha:C\simeq C'$; the choice between $W^1$ and $-W^1$ is the sign
$\varepsilon$.  The equality of Abel maps then has the form
$f'\alpha=\varepsilon\beta f+c$ for a translation $c$.  The incidence and
linear-system descriptions show that this data is unique in the nonhyperelliptic
case, while in the hyperelliptic case the two signs give the two possible
choices.
\end{proof}


\begin{corollary}[Torelli over a perfect field]
\label{cor:milne-torelli-perfect-source}

\dcref{MI:ch3:12.2}
If two genus-$\ge2$ curves over a perfect field have isomorphic canonically
polarized Jacobians over that field, then the curves are isomorphic over it.
\end{corollary}
\begin{proof}
After base change to an algebraic closure, Torelli gives the curve
isomorphism.  Its uniqueness makes it Galois invariant, and perfectness gives
effective descent.
\end{proof}

\begin{corollary}[Torelli finiteness]
\label{cor:milne-torelli-finiteness-source}

\dcref{MI:ch3:12.3}
For fixed $k,g,S$, the polarized-Jacobian map injects curves of genus $g\ge2$
with good reduction outside $S$ into principally polarized abelian varieties
with good reduction outside $S$.
\end{corollary}
\begin{proof}
Good reduction of a curve gives good reduction of its canonically polarized
Jacobian, and Torelli over the base field proves injectivity.
\end{proof}

\begin{corollary}[Torelli and Mordell]
\label{cor:milne-torelli-mordell-source}

\dcref{MI:ch3:12.4}
\uses{cor:milne-torelli-finiteness-source,thm:milne-shafarevich-implies-mordell-source}
Finiteness of principally polarized abelian varieties of fixed dimension and
good reduction outside $S$ implies Mordell's conjecture.
\end{corollary}
\begin{proof}
Torelli converts the assumed polarized-abelian finiteness into Shafarevich
finiteness for curves.  The theorem that Shafarevich implies Mordell then
applies.
\end{proof}

\section{Torelli proof: incidence lemmas}
\begin{lemma}[Canonical theta symmetry]
\label{lem:milne-torelli-theta-symmetry-source}

\dcref{MI:ch3:13.1}
For $a\in J(k)$, $(W_a^{g-1})^*=W_{-a}^{g-1}$.
\end{lemma}

\begin{proof}

Riemann--Roch gives an effective complement to every divisor of degree $g-1$
in a canonical divisor; applying the Abel map and then replacing $a$ by $-a$
gives both inclusions.

\end{proof}


\begin{lemma}[Theta incidence criterion]
\label{lem:milne-torelli-incidence-source}

\dcref{MI:ch3:13.2}
For $0\le r\le g-1$, $W_a^r\subset W_b^{g-1}$ iff
$a\in W_b^{g-1-r}$.
\end{lemma}

\begin{proof}

One implication adds effective divisors.  For the converse, use Riemann--Roch
to complete a degree-$g-1$ divisor to a canonical divisor and compare its Abel
class; this produces the required effective divisor of degree $g-1-r$.

\end{proof}


\begin{lemma}[Theta intersections]
\label{lem:milne-torelli-intersections-source}

\dcref{MI:ch3:13.3}
For $0\le r\le g-1$,
\[
W^{g-1-r}=\bigcap_{a\in W^r}W_a^{g-1},\qquad
(W^{g-1-r})^*=\bigcap_{a\in W^r}(W_a^{g-1})^*.
\]
\end{lemma}

\begin{proof}

Apply the preceding incidence criterion in both directions and use the
canonical-theta symmetry.

\end{proof}


\begin{lemma}[Generic theta intersection decomposition]
\label{lem:milne-torelli-generic-intersection-source}

\dcref{MI:ch3:13.4}
If $a+x=b+y$ with $x\in W^1$, $y\in W^{g-1-r}$ and
$W_a^{r+1}\not\subset W_b^{g-1}$, then
$W_a^{r+1}\cap W_b^{g-1}=W_{a+x}^r\cup S$, where
$S=W_a^{r+1}\cap(W_{y-a}^{g-2})^*$.
\end{lemma}

\begin{proof}

Write $x=f(X)$ and $y=f(Y)$.  A point in the intersection gives effective
$D,D'$ with $D+Y\sim D'+X$.  If the divisors coincide, the point lies in
$W_{a+x}^r$; otherwise Riemann--Roch puts it in the second component.  The
reverse inclusion is immediate from the definitions.

\end{proof}


\begin{lemma}[Theta divisor recovers the Abel curve]
\label{lem:milne-torelli-recovery-source}

\dcref{MI:ch3:13.5}
If $a\in J(k)$ and $W^1\not\subset W_a^{g-1}$, then the intersection is the
effective divisor $D(a)$ of degree $g$ with $f(D(a))=a+\kappa$.
\end{lemma}
\begin{proof}[Source proof]
The incidence criterion identifies $W^1\cap W_a^{g-1}$ with the effective
divisor whose Abel class is $a+\kappa$.  Successively applying the generic
intersection decomposition removes the extraneous components and recovers a
translate of $W^1$ from the polarized theta incidence data.
\end{proof}
